\section{Conclusions}
\label{sec:conclusions}

We presented the first method for both camera calibration and pose estimation from multiple images of known 3D cylinders using homotopy continuation.
By revisiting the geometric constraints of \cite{gummeson22pose}, we were able to  formulate these two tasks within a unified framework flexible enough to adapt to different scenarios on the fly. By exploiting custom silouette-based homotopy continuation solver, a tailored initialization strategy, and a novel correction step we are able to achieve robust and accurate solutions compared to previous works. 

\paragraph{Main contributions} of our work is a flexible homotopy continuation framework to solve high degree polynomial systems arising from geometric computer vision problems.
Specifically, in this paper, we detailed the steps to use the method for camera calibration and pose estimation from multiple images of known 3D cylinders, in different settings and configurations.

An initialization strategy based on a random walk of the solution space, enforcing subsequent iterations to be different enough, the solution space to be traversed evenly, and a priori knowledge to be incorporated efficiently.

A homotopy function specifically designed for the cylinder case to be scale independent by making use of angles, and maximise path convergence.

A correction step general enough to be applied to any method performing both camera calibration and pose estimation.

A comprehensive list of pratical applications of the method in real-world scenarios like camera calibration using a calibration rig, automatic periodic camera calibration, and pose tracking of a video.

\paragraph{Experimental results} demonstrate superior performance of our approach with respect  to existing solutions, both in terms of i) \emph{generality},  ii) \emph{accuracy} and iii) \emph{robustness}.

That said preliminary steps for the method like edge detection and line-cylinder matching pose a hard challenge, especially in very lit or very dim environments, where the silhouette lines and the shadows generated by the cylinders get confused frequently.
Detecting lines using Canny edge detector~\cite{Canny1986}, RANSAC~\cite{USAC} or Hough transform~\cite{Yousefi2018} suffer harsh light conditions.

Matching based on parallelism by testing angle between them works well for the studied scenes but does not work in scenarios where the angle between the camera and the observed cylinders is very steep, or other hard surfaces edges mostly align with the cylinder edges. That said this method is complementary to classic point calibration. Our method thrives in simple scenes, while point-based methods thrive in an abundancy of landmarks.

Another initially not considered challenge is the effect of even small errors on short lines, even a single pixel difference near the end points can skew the line angle significantly leading to large errors in the reconstruction.

\paragraph{Future work}.

Our work opens up several exiting research directions.

First of all, we plan to study an automatic matching criteria for silhouette line-cylinder pairs. The current matching method is based on the angle between the lines, which works well in the studied scenes, but does not generalize well to more complex scenarios. We plan to study automatic matching criteria for silhouette line-cylinder pairs, which would allow our method to be applied to a wider range of scenes. Right now the most promising techniques are Holistically-Nested Edge Detection~\cite{Xie2015} and Liu work on surfaces of revolution~\cite{liu2019realtime}. The second one particularly leverages the symmetry of the cylinder under different rotations,
and camera poses

Second, we plan to apply our method to simplifier scenarios in which our general method can strive and leverage additional constraint to improve speed and accuracy. Right now the method as is is designed to work with a flexible amount of cylinders, a flexible amount of images, a flexible amount of unknown camera intrinsics. This flexibility comes at the cost of computational complexity, which can be reduced by simplifying the problem. Having to deal with a variable number of cylinders in any position also makes edge detection, matching, and path tracking more challenging.

Third, extending the previous point, we don't merely want to simplify the problem, but to find new constraints arising from special cases of the one treated in this paper.

Fourth our custom homotopy function demonstrated to be effective in the cylinder case, but is just the starting point in a growing field. Homotopy continuation is a powerful tool that recently gained popularity in computer vision, and we plan to come up with a vast range of homotopy functions that try to be geometrically aware of the problem at hand.

Fifth we are confident that the very same cylinder constraints can be adapted to other rich geometric primitives, aside from the obvious quadrics, such as spheres, cones, and ellipsoids. For example we are curious wether lines can be treated as cylinders of infinitely small radius. This would allow us to apply our method to a wider range of scenes and problems.

Sixth, homotopy continuation has mostly thrived in CPU predominant fields but computer vision heavily relies on GPUs. We plan to study how to adapt our method to run on GPUs, which would allow us to leverage the massive parallelism of modern GPUs and speed up path tracking by sharing gradients and exploring multiple paths simultaneously.

Finally we plan to study how to apply our method to practical computer vision problems, such as 3D reconstruction, visual odometry, and SLAM. We made a quick overview of the direct applications but it was merely a showcase for tasks that on their own are worth a research paper. The flexibility of our method makes it a promising candidate for these tasks, and we are excited to explore its potential in these areas.
